\subsection{Definitions and evaluation metrics}
\label{sec:metrics}

This section fixes the quantities the rest of the document reports. It exists
because several of the mistakes recorded here were mistakes of \emph{metric}
and not of model: a figure that was correct as computed and wrong as
interpreted. Stating the definitions precisely is what made those visible.

\subsubsection{Matching predictions to ground truth}

For two axis-aligned boxes $a, b$ the intersection over union is
\begin{equation}
  \iou(a,b) \;=\; \frac{|a \cap b|}{|a \cup b|} \;\in\; [0,1].
  \label{eq:iou}
\end{equation}

Given predictions $\hat{B} = \{\hat{b}_1, \dots, \hat{b}_m\}$ sorted by
descending confidence and ground truth $B = \{b_1, \dots, b_n\}$, each
prediction is matched greedily to the highest-$\iou$ unmatched ground-truth
box exceeding a threshold $\tau$. Writing $\mathrm{TP}$, $\mathrm{FP}$ and
$\mathrm{FN}$ for the resulting counts,
\begin{equation}
  \precision = \frac{\mathrm{TP}}{\mathrm{TP} + \mathrm{FP}}, \qquad
  \recall = \frac{\mathrm{TP}}{\mathrm{TP} + \mathrm{FN}}.
  \label{eq:pr}
\end{equation}

Throughout this document $\tau = 0.3$. That is looser than the $0.5$
convention, and the choice is deliberate: the operational question is whether
a person or a plume was \emph{found}, not whether its box was drawn tightly.
A stricter threshold answers a localisation question this system does not ask.

\subsubsection{Box-level and image-level evaluation}
\label{sec:twolevels}

\emph{Box-level} metrics use \cref{eq:pr} over instances. \emph{Image-level}
metrics reduce each image to a binary decision — did the detector fire at all
— and compute the same quantities over images. The image-level false-alarm
rate over a negative set $N$ (images containing no instance of the hazard) is
\begin{equation}
  \mathrm{FAR} \;=\; \frac{|\{\, i \in N : \hat{y}_i = 1 \,\}|}{|N|}.
  \label{eq:far}
\end{equation}

Both are reported everywhere, because either alone conceals a failure the
other exposes.

\begin{recorded}{The failure that image-level metrics hid}
The people-in-water detector's exported head has six logit slots for five
classes. A leading no-object slot was assumed, and every label shifted by one:
boats were reported as swimmers at $0.84$ confidence while real swimmers were
discarded. It nearly survived review because image-level precision still read
$0.83$ — boats and swimmers co-occur in the same frames, so an image with a
boat called ``swimmer'' is still an image containing a swimmer. Only
box-level $\iou$ matching exposed it, at a recall of $0.002$.
\end{recorded}

\subsubsection{Recall-weighted \texorpdfstring{$F$}{F} scores}

The harmonic family
\begin{equation}
  F_\beta \;=\; \frac{(1 + \beta^{2})\, \precision \recall}
                     {\beta^{2} \precision + \recall}
  \label{eq:fbeta}
\end{equation}
is reported at $\beta = 2$ as the headline, a convention borrowed from the
RipVIS benchmark~\citep{ripvis2025}. The justification is not stylistic. At
the symmetric point $\precision = \recall = p$,
\begin{equation}
  \frac{\partial F_\beta / \partial \recall}
       {\partial F_\beta / \partial \precision}
  \;=\; \beta^{2},
  \label{eq:fbeta-ratio}
\end{equation}
so $\beta = 2$ states explicitly that a point of recall is worth four points
of precision here. That is the correct trade for this system for an asymmetric
reason: the validator downstream can discard a false alarm, but nothing
downstream can recover a person who was never detected.

\subsubsection{Segmentation metrics and pooling}
\label{sec:pooling}

For binary masks the water $\iou$ on a single image is \cref{eq:iou} over
pixel sets. Over a corpus of $n$ images there are two ways to aggregate, and
they answer different questions:
\begin{align}
  \iou_{\text{pooled}} &= \frac{\sum_{i=1}^{n} |P_i \cap G_i|}
                               {\sum_{i=1}^{n} |P_i \cup G_i|},
  \label{eq:iou-pooled} \\[2pt]
  \iou_{\text{strat}}  &= \frac{\sum_{i \in S} |P_i \cap G_i|}
                               {\sum_{i \in S} |P_i \cup G_i|},
  \qquad S \subset \{1,\dots,n\}.
  \label{eq:iou-strat}
\end{align}

\Cref{eq:iou-pooled} weights every image by the size of its union, so a corpus
dominated by easy cases is dominated by easy cases in the score. When the
subset $S$ that the model exists to handle is a small minority, the pooled
figure flatters it, and the more severe the case the smaller its influence.

This is not hypothetical. FloodNet's~\citep{rahnemoonfar2021floodnet} labelled
release contains \AerialLabelled{} pairs of which only \AerialFlooded{} are
flooded — \AerialFloodedShare\,\% of the corpus. The aerial flood segmenter
scores $\iou_{\text{pooled}} = \AerialOverallIoU$ but
$\iou_{\text{strat}} = \AerialFloodedIoU$ on the flooded subset. Model
selection in this project uses the stratified figure, and \cref{sec:flood}
publishes it as the headline with the pooled one beside it, clearly labelled.
Selecting on the pooled figure would have chosen a model that is better at
recognising dry ground.

\subsubsection{Interval estimation}
\label{sec:intervals}

Half of this project's figures rest on small samples: \BuildingTestImages{}
held-out images for building access, \FireNarrowNegatives{} negatives in the
first fire baseline. A bare point estimate from such a sample invites a reader
to believe it to three decimal places, so every proportion in this document
whose denominator is known is reported with a 95\,\% \emph{Wilson score}
interval~\citep{wilson1927probable}. For $k$ successes in $n$ trials with
$\hat{p} = k/n$ and $z = 1.96$,
\begin{equation}
  \mathrm{CI}_{95} \;=\;
  \frac{\hat{p} + \dfrac{z^{2}}{2n}}{1 + \dfrac{z^{2}}{n}}
  \;\pm\;
  \frac{z\sqrt{\dfrac{\hat{p}(1-\hat{p})}{n} + \dfrac{z^{2}}{4n^{2}}}}
       {1 + \dfrac{z^{2}}{n}}.
  \label{eq:wilson}
\end{equation}

Wilson over the textbook normal approximation
$\hat{p} \pm z\sqrt{\hat{p}(1-\hat{p})/n}$, for a reason this project's data
makes concrete instead of academic~\citep{brown2001interval}. At
$\hat{p} = 1$ the normal interval collapses to a point: building access
detects a person in $12$ of $12$ positive held-out images, and the normal
interval reports $[1.000, 1.000]$ — perfect certainty inferred from twelve
observations. \Cref{eq:wilson}, derived by inverting the score test rather
than by assuming normality of the estimate, gives
\BuildingImageRecallCI. That interval is the honest statement, and it is what
\cref{sec:building} reports.

\Cref{tab:intervals} collects the cases where the interval changes how a
figure should be read.

\begin{table}[htbp]
\centering\small
\caption{Point estimates with the 95\,\% Wilson interval their sample size
supports. The last column is the reading the interval permits, which in three
of these rows is materially weaker than the point estimate alone suggests.
Computed from the measurement files.}
\label{tab:intervals}
\begin{tabularx}{\textwidth}{l S[table-format=1.4] c L}
\toprule
Quantity & {Estimate} & 95\,\% interval & What it supports \\
\midrule
Fire image-level FAR, \FireNegatives{} negatives
  & \FireImageFAR & \FireImageFARCI
  & Precise enough to quote. The motivating claim is safe. \\
Fire image-level FAR, \FireNarrowNegatives{} negatives
  & \FireNarrowFAR & \FireNarrowFARCI
  & Consistent with anything from two-thirds to nearly all. Withdrawn. \\
Swimmer box recall, \SwimmerInstances{} instances
  & \SwimmerBoxRecall & \SwimmerBoxRecallCI
  & Tight. The strongest evidence in the project. \\
Civilian box recall, \BuildingCivilians{} instances
  & \BuildingCivilianRecall & \BuildingCivilianRecallCI
  & Solid at instance level, from \BuildingTestImages{} images. \\
Building image-level recall, 12 images
  & \BuildingImageRecall & \BuildingImageRecallCI
  & Cannot exclude one image in four being missed. \\
Building image-level FAR, 12 negatives
  & \BuildingImageFAR & \BuildingImageFARCI
  & Almost uninformative. Reported for completeness only. \\
Scene router accuracy, \RouterHeldOut{} images
  & \RouterAccuracy & \RouterAccuracyCI
  & Tight, but see \cref{sec:routing} on task difficulty. \\
\bottomrule
\end{tabularx}
\end{table}

\subsubsection{Class merging}

Two detectors emit more classes than VIGÍA consumes. Merging is a partial
surjection
\begin{equation}
  \mu : \mathcal{K}_{\text{model}} \longrightarrow
        \mathcal{K}_{\text{system}} \cup \{\bot\},
  \label{eq:merge}
\end{equation}
where $\bot$ discards a class. For building access $|\mathcal{K}_{\text{model}}| = 28$
and $|\mathcal{K}_{\text{system}}| = 5$; for people in water, five model classes
reduce to one consumed class. Crucially, $\mu$ is applied at the detector
boundary and only a designated subset
$\mathcal{A} \subseteq \mathcal{K}_{\text{system}}$ — the \emph{alarm} classes
— reaches the validator. Everything else is operator context. A rescue team on
site is the expected state at an incident already being handled, and a boat is
not an emergency; routing either to the validator would alarm on every rescue
worker and every vessel in frame.
